\subsection{THE MULTIPLICATIVE GROUP C*}

We know from Chapter III, § 6, no. 7 that the topology induced on the multiplicative group \(\mathbf{C}^*\) of non-zero complex numbers is compatible with the group structure of \(\mathbf{C}^*\) . Since \(\mathbf{C}^*\) is open in \(\mathbf{C}\) , it follows that \(\mathbf{C}^*\) is a locally compact topological group (Chapter I, § 9, no. 7, Proposition 13) and therefore complete (with respect to the multiplicative uniformity; cf.~Chapter III, § 6, no. 8, Proposition 8). The multiplicative group \(\mathbf{R}_+^*\) of real numbers \(>0\) is a closed subgroup of \(\mathbf{C}^*\) . Another subgroup is the set \(\mathbf{U}\) of complex numbers of absolute value 1, which is identified with the unit circle \(\mathbf{S}_1\) of \(\mathbf{R}^2\) , and is therefore a compact group. Moreover:

PROPOSITION 1. The topological group \(\mathbf{C}^*\) is isomorphic to the product of the topological groups \(\mathbf{R}^*\) and \(\mathbf{U}\) .

For the mapping \(z \to \left( |z|, \frac{z}{|z|} \right)\) is a homeomorphism of \(\mathbf{C}^*\) onto \(\mathbf{R}_+^* \times \mathbf{U}\) (Chapter VI, § 2, no. 3, Proposition 3); and it follows immediately that it is an isomorphism of the group structures.

The topological group \(\mathbf{R}_{+}^{*}\) is already known to be isomorphic to the additive group \(\mathbf{R}\) (Chapter V, § 4, Theorem 1); the study of the topological group \(\mathbf{C}^{*}\) is therefore reduced to that of \(\mathbf{U}\) , which we shall consider in § 2.

